\section{Embodied calibration and the proposed updating loop}
\subsection{Measurement before interpretation}
The watch sequence proposes a particularly clear methodological rule: estimate a bodily variable before consulting a measurement \srcs{64--69}. Its value is that it makes disagreement observable. It does not establish that internal sensing outperforms instruments. A wearable has sampling, artifact, and algorithmic error; a subjective estimate has perceptual, memory, expectation, and reporting error. Neither should silently define the other as correct.

Interoceptive accuracy, self-reported sensibility, and metacognitive awareness are distinguishable constructs in the relevant literature \cite{garfinkel}. Pulse estimation, respiration estimation, fatigue ratings, muscular tension, and recovery scores also require different measurement models. A proprietary ``coherence'' score has no usable scientific meaning here until its algorithm, units, validation, and intended interpretation are specified. Missing watch data during an exercise remain missing; recollection is not a replacement physiological recording.

For participant $i$ and trial $t$, record an estimate $\hat y_{it}$, confidence, and the exact time window before revealing a synchronized reference value $y^{\rm ref}_{it}$. Then calculate
\begin{align}
 e_{it}&=\hat y_{it}-y^{\rm ref}_{it},\\
 \operatorname{MAE}_i&=\frac{1}{n_i}\sum_t|e_{it}|,
 &\operatorname{Bias}_i&=\frac{1}{n_i}\sum_t e_{it}.
\end{align}
An uncertainty interval $I_{it}$ should additionally be assessed for coverage and width. Always reporting a very wide interval is not useful calibration. Error is relative to the reference procedure, whose own uncertainty and invalid-trial criteria must be reported. Do not pool heart-rate errors in beats per minute with respiration errors in breaths per minute without a justified normalization.

An estimate can be accurate because the person knows a typical resting heart rate or infers elapsed time. Heartbeat counting therefore cannot alone identify interoceptive perception; experimental evidence demonstrates substantial non-interoceptive contributions \cite{desmedt}. A stronger battery combines multiple tasks, randomized windows, recorded strategies, baseline-knowledge controls, and a prespecified reference modality. Attending to respiration may also change it, so the estimation and reference windows must coincide.

\subsection{Calibration, autonomy, and sensor use}
The narrative rejects both replacing a device with unexplained teacher authority and refusing useful technology to maintain an identity \srcs{68--69}. A real protocol needs explicit criteria: task-specific error tolerances, uncertainty coverage, repeated independent trials, transfer to a new setting, and a plan for detecting drift. Tolerances should follow the intended use and measurement precision, not be invented from the story. Human estimates must not replace monitoring required for care or safety.

Feedback can be compared with delayed feedback or a matched educational control. Removing a display tests performance without visible feedback; it does not necessarily require removing the sensor. Keeping the reference recording while hiding the display avoids creating an unnecessary data gap. Calibration training is a hypothesis about learning, not a consequence of voluntary device surrender.

\subsection{Embodied updating as a decomposable intervention}
The source combines attention to sensation, an EFT-like tapping sequence, breathing, a verbal revision cue, and reassessment \srcs{71--74}. The most defensible scientific translation is a multicomponent experimental package, provisionally named \emph{embodied updating}. The sequence is: record the current experience and belief; apply the assigned attention/sensory routine; reassess; identify any new evidence; state whether the belief, affect, or intended action changed. The package is not presented here as a validated therapy.

The word ``updating'' can describe at least three different events: a change in reported affect, a change in a probabilistic belief about an external proposition, or a change in intended behavior. These outcomes must not be conflated. Feeling calmer does not automatically make a prediction more accurate, and feeling more anxious does not automatically make an exercise unsuccessful. A legitimate result may be greater awareness of fatigue with no change in an external belief.

Breathing research provides a relevant but bounded connection: one randomized study compared structured breathing practices with mindfulness meditation and measured mood and physiological outcomes \cite{balban}. It did not test this narrative's complete routine, factual belief accuracy, or meridian mechanisms. Component studies of EFT likewise do not settle which aspect of a new package would be active \cite{waite,church}. Those studies differ in populations, controls, and tasks; they cannot be pooled informally into proof that the present sequence works. Sensory anchoring, expectancy, rhythmic attention, exposure, task interruption, social support, and the evidence prompt are competing explanations.

No claim is made that the listed tapping points possess special energetic properties. To test location specificity requires a matched alternative-location comparison holding touch, rhythm, expectancy, and instruction as constant as practicable. To test breath pacing requires a matched comparator and a defined comfortable protocol. To test cognitive framing requires a neutral verbal control. The factorial design in Section~\ref{sec:protocols} targets these separable questions.

\subsection{Synchrony and interpersonal inference}
The paired watch response during fictional travel is an example of underdetermination \srcs{26--27}. Shared stimuli, common motion artifacts, synchronized instructions, social interaction, chance, and selective inspection can all produce apparent synchrony. A future ordinary physiological study would need synchronized clocks, raw signal quality checks, predefined windows, baseline and shuffled-pair controls, and models accounting for autocorrelation. Selecting the largest cross-correlation after trying many lags inflates apparent evidence.

The same caution applies to behavioral cues. A gaze shift, reduced blinking, or muscle tension is not an individually validated lie detector. The deception-cue literature finds many weak or absent associations \cite{depaulo}. This does not imply that all nonverbal information is useless; it means that a scientific claim needs predictive performance and error rates in its intended setting. Consent to observe, collect, infer, or share personal data remains distinct from the ability to notice a cue.
